\documentclass[a4paper,landscape,8pt]{extarticle}
\usepackage[landscape,margin=0.7cm,top=0.5cm,bottom=0.7cm]{geometry}
\usepackage[utf8]{inputenc}
\usepackage[T1]{fontenc}
\usepackage[ngerman]{babel}
\usepackage{helvet}
\renewcommand{\familydefault}{\sfdefault}
\usepackage{xcolor}
\usepackage{tikz}
\usetikzlibrary{positioning}
\usepackage{enumitem}
\usepackage{multicol}
\usepackage{array}
\usepackage{fancyhdr}
\usepackage{amssymb}

% Farben
\definecolor{mailblue}{HTML}{1565C0}
\definecolor{textgreen}{HTML}{2E7D32}
\definecolor{attachorange}{HTML}{E65100}
\definecolor{lightgray}{HTML}{F5F5F5}
\definecolor{bordergray}{HTML}{BBBBBB}

\pagestyle{fancy}
\fancyhf{}
\renewcommand{\headrulewidth}{0pt}
\fancyfoot[C]{\footnotesize Seite \thepage{} von 3}

\setlength{\parindent}{0pt}
\setlength{\parskip}{2pt}

% Kompakte Listen
\setlist{nosep, leftmargin=1em, itemsep=0pt, topsep=1pt}

\begin{document}

% ============================================================================
% SEITE 1: INFORMATIONSBLATT
% ============================================================================

\noindent{\Large\bfseries E-Mails schreiben: Der digitale Brief} \hfill
{\footnotesize\textbf{Lernziele:} Mail verfassen | Anrede wählen | Foto anhängen}

\vspace{1mm}
\hrule
\vspace{1mm}

% Drei-Spalten-Layout
\begin{minipage}[t]{0.32\linewidth}
\begin{center}
\colorbox{mailblue!10}{%
\begin{minipage}{0.95\linewidth}
\vspace{2mm}
\begin{center}
{\Large\bfseries\textcolor{mailblue}{1. Eine neue Mail}}\\[1pt]
{\footnotesize\itshape Der Aufbau}
\end{center}
\vspace{2mm}
\end{minipage}%
}
\end{center}

\vspace{1mm}

\textbf{So startest du:}
\begin{enumerate}
\item Mail-App öffnen
\item Quadrat mit Stift antippen (unten rechts)
\item Neues Fenster öffnet sich
\end{enumerate}

\textbf{Die Felder:}
\begin{itemize}
\item \textbf{An:} E-Mail-Adresse des Empfängers
\item \textbf{CC:} Kopie an weitere Personen (alle sehen es)
\item \textbf{Betreff:} Kurze Überschrift -- worum geht es?
\item \textbf{Text:} Deine Nachricht
\end{itemize}

\textbf{Analogie:} Wie beim Brief:
\begin{itemize}
\item An = Adresse auf dem Umschlag
\item Betreff = Überschrift
\item Text = Der eigentliche Brief
\end{itemize}

\end{minipage}%
\hfill%
\begin{minipage}[t]{0.32\linewidth}
\begin{center}
\colorbox{textgreen!10}{%
\begin{minipage}{0.95\linewidth}
\vspace{2mm}
\begin{center}
{\Large\bfseries\textcolor{textgreen}{2. Anrede \& Grußformel}}\\[1pt]
{\footnotesize\itshape Der gute Ton}
\end{center}
\vspace{2mm}
\end{minipage}%
}
\end{center}

\vspace{1mm}

\textbf{Formelle Anrede:}\\
(Behörden, Firmen, Unbekannte)
\begin{itemize}
\item ``Sehr geehrte Frau Müller,''
\item ``Sehr geehrter Herr Schmidt,''
\item ``Sehr geehrte Damen und Herren,''
\end{itemize}

\textbf{Informelle Anrede:}\\
(Familie, Freunde)
\begin{itemize}
\item ``Liebe Maria,'' / ``Lieber Peter,''
\item ``Hallo Maria,''
\end{itemize}

\textbf{Grußformeln:}
\begin{itemize}
\item \textbf{Formell:} ``Mit freundlichen Grüßen''
\item \textbf{Informell:} ``Liebe Grüße'' oder ``Viele Grüße''
\end{itemize}

\vspace{2mm}
\fcolorbox{textgreen}{textgreen!5}{%
\begin{minipage}{0.92\linewidth}
\footnotesize
\textbf{Tipp:} Nach der Anrede kommt ein \textbf{Komma}, dann eine Leerzeile, dann der Text.
\end{minipage}%
}

\end{minipage}%
\hfill%
\begin{minipage}[t]{0.32\linewidth}
\begin{center}
\colorbox{attachorange!10}{%
\begin{minipage}{0.95\linewidth}
\vspace{2mm}
\begin{center}
{\Large\bfseries\textcolor{attachorange}{3. Foto anhängen}}\\[1pt]
{\footnotesize\itshape Dateien mitschicken}
\end{center}
\vspace{2mm}
\end{minipage}%
}
\end{center}

\vspace{1mm}

\textbf{Methode 1:}
\begin{enumerate}
\item Im Text \textbf{tippen und halten}
\item ``Foto oder Video einfügen'' wählen
\item Foto aus der Galerie aussuchen
\end{enumerate}

\textbf{Methode 2:}
\begin{enumerate}
\item Über der Tastatur: \textbf{$<$} Symbol
\item Dann Büroklammer oder Kamera
\item ``Fotomediathek'' wählen
\end{enumerate}

\textbf{Wichtig:}
\begin{itemize}
\item Fotos können groß sein (mehrere MB)
\item iPhone fragt nach Größe
\item ``Groß'' oder ``Tatsächlich'' wählen
\end{itemize}

\textbf{Senden:}
\begin{itemize}
\item Alles fertig?
\item \textbf{Blauer Pfeil} oben rechts
\item Mail ist unterwegs!
\end{itemize}

\end{minipage}

\vspace{1mm}

% Zusammenfassung
\begin{center}
\fcolorbox{bordergray}{lightgray}{%
\begin{minipage}{0.98\linewidth}
\centering
\vspace{1.5mm}
\textbf{Zusammenfassung:}
\textcolor{mailblue}{\textbf{Neue Mail}} = Quadrat mit Stift ~$|$~
\textcolor{textgreen}{\textbf{Betreff}} nicht vergessen! ~$|$~
\textcolor{attachorange}{\textbf{Foto}} = tippen \& halten oder $<$ Symbol ~$|$~
\textbf{Senden} = blauer Pfeil
\vspace{1.5mm}
\end{minipage}%
}
\end{center}

\newpage

% ============================================================================
% SEITE 2: ÜBUNGEN
% ============================================================================

\begin{center}
{\LARGE\bfseries Übungen: E-Mails schreiben}\\[3pt]
{\normalsize Schau dir Seite 1 an und beantworte die Fragen}
\end{center}

\vspace{3mm}
\hrule
\vspace{4mm}

\begin{multicols}{2}

\textbf{\large Teil A: Die richtige Anrede}

Welche Anrede passt? Schreibe sie auf.

\vspace{2mm}

\renewcommand{\arraystretch}{2}
\begin{tabular}{@{}p{5.5cm}p{6cm}@{}}
1. An deine beste Freundin Maria: & \dotfill \\
2. An deinen Arzt Dr. Schmidt: & \dotfill \\
3. An eine Firma, Ansprechpartner unbekannt: & \dotfill \\
4. An deinen Enkel Tom: & \dotfill \\
\end{tabular}
\renewcommand{\arraystretch}{1}

\vspace{6mm}

\textbf{\large Teil B: Gute Betreffzeilen}

Welcher Betreff ist besser? Kreuze an.

\vspace{2mm}

\renewcommand{\arraystretch}{1.8}
\begin{tabular}{@{}p{5cm}p{5cm}c@{}}
1a) ``Hallo'' & 1b) ``Terminanfrage Zahnarzt'' & $\square$ a ~$\square$ b \\
2a) ``Frage'' & 2b) ``Frage zu meiner Rechnung'' & $\square$ a ~$\square$ b \\
3a) ``Fotos vom Geburtstag'' & 3b) ``Fotos'' & $\square$ a ~$\square$ b \\
4a) ``Wichtig!!!'' & 4b) ``Adressänderung mitteilen'' & $\square$ a ~$\square$ b \\
\end{tabular}
\renewcommand{\arraystretch}{1}

\vspace{6mm}

\textbf{\large Teil C: Reihenfolge}

Nummeriere die Schritte beim E-Mail-Schreiben (1--6):

\vspace{2mm}

\renewcommand{\arraystretch}{1.8}
\begin{tabular}{@{}cp{8cm}@{}}
$\square$ & Grußformel und Namen schreiben \\
$\square$ & Auf ``Senden'' (blauer Pfeil) tippen \\
$\square$ & E-Mail-Adresse bei ``An:'' eingeben \\
$\square$ & Mail-App öffnen \\
$\square$ & Text schreiben \\
$\square$ & Betreff eingeben \\
\end{tabular}
\renewcommand{\arraystretch}{1}

\columnbreak

\textbf{\large Teil D: Praktisch üben}

Schreibe eine kurze E-Mail auf. Fülle die Felder aus:

\vspace{3mm}

\textbf{Situation:} Du schreibst an deine Tochter und möchtest sie nächste Woche besuchen.

\vspace{3mm}

\fcolorbox{bordergray}{white}{%
\begin{minipage}{0.95\linewidth}
\vspace{3mm}
\textbf{An:} \dotfill \\[4mm]
\textbf{Betreff:} \dotfill \\[4mm]
\hrule
\vspace{3mm}
\textit{(Anrede)}\\[6mm]
\rule{\linewidth}{0.4pt}\\[6mm]
\rule{\linewidth}{0.4pt}\\[6mm]
\rule{\linewidth}{0.4pt}\\[6mm]
\textit{(Grußformel + Name)}\\[6mm]
\rule{\linewidth}{0.4pt}
\vspace{3mm}
\end{minipage}%
}

\vspace{5mm}

\textbf{Bonus:} Wie würdest du ein Foto anhängen?

\vspace{4mm}
\rule{\linewidth}{0.4pt}

\end{multicols}

\newpage

% ============================================================================
% SEITE 3: CHECKLISTE
% ============================================================================

\begin{center}
{\LARGE\bfseries Checkliste: Das habe ich verstanden}\\[3pt]
{\normalsize Geh die Liste durch und hake ab, was du sicher weißt}
\end{center}

\vspace{3mm}
\hrule
\vspace{6mm}

% Drei Boxen nebeneinander
\begin{minipage}[t]{0.31\linewidth}
\begin{center}
\fcolorbox{mailblue}{mailblue!8}{%
\begin{minipage}{0.92\linewidth}
\vspace{3mm}
\begin{center}
{\large\bfseries\textcolor{mailblue}{Neue Mail erstellen}}
\end{center}
\vspace{2mm}
\begin{itemize}[label=$\square$, itemsep=5pt]
\item Ich finde das \textbf{``Neue Mail''}-Symbol
\item Ich weiß, wo ich die \textbf{Empfänger-Adresse} eingebe
\item Ich vergesse den \textbf{Betreff} nicht
\item Ich weiß, wo der \textbf{Senden}-Button ist
\end{itemize}
\vspace{3mm}
\end{minipage}%
}
\end{center}
\end{minipage}%
\hfill%
\begin{minipage}[t]{0.31\linewidth}
\begin{center}
\fcolorbox{textgreen}{textgreen!8}{%
\begin{minipage}{0.92\linewidth}
\vspace{3mm}
\begin{center}
{\large\bfseries\textcolor{textgreen}{Anrede \& Grußformel}}
\end{center}
\vspace{2mm}
\begin{itemize}[label=$\square$, itemsep=5pt]
\item Ich kenne \textbf{formelle} Anreden
\item Ich kenne \textbf{informelle} Anreden
\item Ich weiß, welche \textbf{Grußformel} passt
\item Ich setze nach der Anrede ein \textbf{Komma}
\end{itemize}
\vspace{3mm}
\end{minipage}%
}
\end{center}
\end{minipage}%
\hfill%
\begin{minipage}[t]{0.31\linewidth}
\begin{center}
\fcolorbox{attachorange}{attachorange!8}{%
\begin{minipage}{0.92\linewidth}
\vspace{3mm}
\begin{center}
{\large\bfseries\textcolor{attachorange}{Anhänge \& Senden}}
\end{center}
\vspace{2mm}
\begin{itemize}[label=$\square$, itemsep=5pt]
\item Ich kann ein \textbf{Foto anhängen}
\item Ich weiß, dass Fotos \textbf{groß} sein können
\item Ich finde den \textbf{Senden}-Button
\item Ich prüfe die Mail \textbf{vor dem Senden}
\end{itemize}
\vspace{3mm}
\end{minipage}%
}
\end{center}
\end{minipage}

\vspace{8mm}

% Gesamtverständnis
\begin{center}
\fcolorbox{black}{lightgray}{%
\begin{minipage}{0.95\linewidth}
\vspace{4mm}
\begin{center}
{\large\bfseries Gesamtverständnis}
\end{center}
\vspace{3mm}
\begin{multicols}{2}
\begin{itemize}[label=$\square$, itemsep=6pt]
\item Ich kann eine \textbf{neue E-Mail} beginnen
\item Ich wähle die passende \textbf{Anrede}
\item Ich schreibe einen \textbf{klaren Betreff}
\item Ich kann ein \textbf{Foto anhängen}
\item Ich kann die Mail \textbf{absenden}
\item Ich fühle mich \textbf{sicher} beim E-Mail-Schreiben
\end{itemize}
\end{multicols}
\vspace{3mm}
\end{minipage}%
}
\end{center}

\vspace{8mm}

% Notfall-Notizen
\begin{center}
\fcolorbox{bordergray}{white}{%
\begin{minipage}{0.95\linewidth}
\vspace{4mm}
\begin{center}
{\large\bfseries Meine Notizen zum E-Mail-Schreiben}\\[3pt]
{\footnotesize (Fülle das hier aus und bewahre es auf!)}
\end{center}
\vspace{4mm}
\renewcommand{\arraystretch}{2}
\begin{tabular}{@{}p{0.45\linewidth}p{0.45\linewidth}@{}}
\textbf{Wichtige E-Mail-Adressen:} & \textbf{Meine Standard-Grußformel:} \\
\dotfill & \dotfill \\[3mm]
\dotfill & \\[3mm]
\dotfill & \\
\end{tabular}
\renewcommand{\arraystretch}{1}
\vspace{4mm}
\end{minipage}%
}
\end{center}

\vfill

\begin{center}
\footnotesize\textit{Stand: \today{} -- Vor dem Senden: Empfänger und Text nochmal prüfen!}
\end{center}

\end{document}
